\documentclass[11pt]{article}
\usepackage[T1]{fontenc}
\usepackage{lmodern}
\usepackage{microtype}
\usepackage{geometry}
\usepackage{amsmath,amssymb,amsthm,mathtools}
\usepackage{hyperref}
\geometry{margin=1in}
\hypersetup{colorlinks=true,linkcolor=blue,urlcolor=blue}

\newtheorem{theorem}{Theorem}
\newtheorem{corollary}{Corollary}
\newtheorem{remark}{Remark}
\DeclareMathOperator{\sgn}{sgn}

\title{An Exact Positive-Wavenumber Minimum\\in a Higher-Derivative Dispersion Relation}
\author{}
\date{26 October 2026}

\begin{document}
\maketitle

\begin{abstract}
We give an exact algebraic criterion for the lower real-frequency branch
of a quadratic higher-derivative dispersion relation to attain its unique
minimum at nonzero wavenumber.  The proof establishes strict convexity in
squared wavenumber, the explicit minimizer, and positivity of both
squared-frequency roots throughout the real-pole interval.  The result is
parameter dependent: an oscillatory-frequency minimum does not remove
higher-derivative ghost or high-wavenumber obstructions.
\end{abstract}

\section{Definitions and main result}
Let $q=|\mathbf{k}|^2\geq 0$ and $z=\omega^2$.  Consider the quadratic
dispersion relation
\begin{equation}\label{eq:dispersion}
 K z^2-(1+a q)z+m^2+q+(K+G+T)q^2=0,
 \qquad a:=2K+G.
\end{equation}
Introduce
\begin{equation}\label{eq:defs}
 B:=4KT-G^2,\qquad
 s:=\sqrt{1-4Km^2},\qquad
 d(q):=s^2+2Gq-Bq^2.
\end{equation}

\begin{theorem}[Exact positive-wavenumber minimum]\label{thm:main}
Assume
\begin{equation}\label{eq:base}
 K>0,\qquad T>0,\qquad B>0,\qquad 0<4Km^2<1,
\end{equation}
and the strict finite-scale condition
\begin{equation}\label{eq:condition}
 \frac{G(1-s)}{2K}>s.
\end{equation}
Then $G>0$ and $a>0$.  The two squared-frequency branches are
\begin{equation}\label{eq:branches}
 z_{\pm}(q)=\frac{1+aq\pm\sqrt{d(q)}}{2K}.
\end{equation}
For nonnegative $q$, the branches are real precisely on
\begin{equation}\label{eq:domain}
 0\leq q\leq q_{\mathrm{coll}},\qquad
 q_{\mathrm{coll}}:=\frac{G+\sqrt{G^2+Bs^2}}{B}.
\end{equation}
The lower branch $z_-$ has a unique global minimum on this interval.
It lies strictly in the interior and is given by
\begin{equation}\label{eq:qstar}
 \boxed{\displaystyle
 q_* = \frac{G-a\sqrt{\frac{G^2+Bs^2}{a^2+B}}}{B}}
 \quad\text{with}\quad
 0<q_*<\frac{G}{B}<q_{\mathrm{coll}}.
\end{equation}
Moreover,
\begin{equation}\label{eq:convex}
 z_-''(q)=\frac{Bs^2+G^2}{2K\,d(q)^{3/2}}>0
 \quad (0<q<q_{\mathrm{coll}}),
\end{equation}
and both squared frequencies are strictly positive on the full
nonnegative real-pole interval.  Thus $|\mathbf{k}_*|=\sqrt{q_*}>0$
is an oscillatory-frequency minimum, not a growing-mode band at that scale.
\end{theorem}

\begin{proof}
\par\medskip\noindent\textbf{Step 1: Compute the discriminant.}\enspace
The discriminant of \eqref{eq:dispersion}, regarded as a quadratic in $z$, is
\begin{align*}
 \Delta(q)
 &=(1+aq)^2-4K\bigl[m^2+q+(K+G+T)q^2\bigr]\\
 &=1-4Km^2+(2a-4K)q
   +\bigl[a^2-4K(K+G+T)\bigr]q^2\\
 &=s^2+2Gq-(4KT-G^2)q^2=d(q),
\end{align*}
where we used $a=2K+G$.  The quadratic formula proves
\eqref{eq:branches}.

\par\medskip\noindent\textbf{Step 2: Determine the real interval.}\enspace
From $0<s<1$ and \eqref{eq:condition}, necessarily $G>0$;
therefore $a=2K+G>0$. Since $B>0$, the two zeros of $d$ are
\[
 q_{\pm}=\frac{G\pm\sqrt{G^2+Bs^2}}{B},
 \qquad q_-<0<q_+=q_{\mathrm{coll}}.
\]
Thus $d(q)>0$ for $0\leq q<q_{\mathrm{coll}}$, $d(q)=0$ at
$q_{\mathrm{coll}}$, and $d(q)<0$ for $q>q_{\mathrm{coll}}$.
This proves the stated real-pole interval.

\par\medskip\noindent\textbf{Step 3: Differentiate.}\enspace
Because $d'(q)=2(G-Bq)$,
\begin{equation}\label{eq:first}
 z_-'(q)=\frac{1}{2K}
 \left(a-\frac{G-Bq}{\sqrt{d(q)}}\right).
\end{equation}
A second differentiation yields
\begin{align*}
 z_-''(q)
 &=\frac{1}{2K}\left[
 \frac{B}{\sqrt{d(q)}}
 +\frac{(G-Bq)^2}{d(q)^{3/2}}\right]\\
 &=\frac{Bd(q)+(G-Bq)^2}{2K\,d(q)^{3/2}}\\
 &=\frac{Bs^2+G^2}{2K\,d(q)^{3/2}}>0.
\end{align*}
Hence $z_-'$ is strictly increasing on the interior real interval.

\par\medskip\noindent\textbf{Step 4: Establish existence and uniqueness.}\enspace
At the origin, using $\sqrt{d(0)}=s$,
\begin{equation}\label{eq:initial}
 z_-'(0)=\frac{a-G/s}{2K}
 =1-\frac{G(1-s)}{2Ks}<0,
\end{equation}
where the last inequality is \eqref{eq:condition}.
Meanwhile $G/B$ lies in $(0,q_{\mathrm{coll}})$ and
\begin{equation}\label{eq:positive}
 z_-'(G/B)=\frac{a}{2K}>0.
\end{equation}
Continuity and the intermediate value theorem give a zero of $z_-'$
in $(0,G/B)$.  Strict positivity of $z_-''$ makes that zero unique.
The derivative is negative before it and positive afterward; hence the
zero is the unique global minimizer on
$[0,q_{\mathrm{coll}}]$ (using continuity at the collision endpoint).

\par\medskip\noindent\textbf{Step 5: Solve the stationary-point equation without extraneous roots.}\enspace
The condition $z_-'(q_*)=0$ is equivalent to
\begin{equation}\label{eq:stationary}
 a\sqrt{d(q_*)}=G-Bq_*.
\end{equation}
Its left side is strictly positive, so $r:=G-Bq_*>0$.
As $q_*=(G-r)/B$,
\[
 d(q_*)=s^2+\frac{G^2-r^2}{B}.
\]
Squaring \eqref{eq:stationary} under the established sign condition
$r>0$ gives
\[
 a^2\left(s^2+\frac{G^2-r^2}{B}\right)=r^2,
 \qquad
 (a^2+B)r^2=a^2(G^2+Bs^2).
\]
We must take the positive root because $r>0$:
\[
 r=a\sqrt{\frac{G^2+Bs^2}{a^2+B}}.
\]
Substitution into $q_*=(G-r)/B$ proves \eqref{eq:qstar}.
The positive-root choice also guarantees that the expression satisfies
the original, unsquared stationary-point equation.

For completeness, since $G,a,B>0$,
\begin{align*}
 q_*>0
 &\iff G>a\sqrt{\frac{G^2+Bs^2}{a^2+B}}\\
 &\iff G^2(a^2+B)>a^2(G^2+Bs^2)\\
 &\iff G>as
 \iff G(1-s)>2Ks,
\end{align*}
which is precisely \eqref{eq:condition}.  Also $r>0$ implies
$q_*<G/B$, and $G/B<q_{\mathrm{coll}}$ follows from
\eqref{eq:domain}.

\par\medskip\noindent\textbf{Step 6: Show the roots are positive.}\enspace
Vieta's identities applied to \eqref{eq:dispersion} yield
\begin{equation}\label{eq:vieta}
 z_+(q)+z_-(q)=\frac{1+aq}{K}>0,
 \qquad
 z_+(q)z_-(q)=
 \frac{m^2+q+(K+G+T)q^2}{K}>0.
\end{equation}
Here $m^2>0$, $K,T,G>0$, and $q\geq0$.
For $0\leq q\leq q_{\mathrm{coll}}$ both roots are real.
Their positive product and positive sum imply that both are strictly
positive.  Accordingly $\omega_-(q)=\sqrt{z_-(q)}$ is real and positive.
Because the square-root map is strictly increasing on $(0,\infty)$,
it attains its unique minimum at the same $q_*$.
Replacing $q$ by $|\mathbf{k}|^2$ proves the positive-wavenumber claim.
\end{proof}

\section{Additional exact consequences}
\begin{corollary}[Exact minimum squared frequency]\label{cor:freq}
Under the assumptions of Theorem~\ref{thm:main},
\begin{equation}\label{eq:zminimum}
 \boxed{\displaystyle
 z_{\min}=z_-(q_*)=
 \frac{B+aG-\sqrt{(a^2+B)(G^2+Bs^2)}}{2KB}>0,}
\end{equation}
and the positive minimum frequency is
$\omega_{\min}=\sqrt{z_{\min}}$.
\end{corollary}
\begin{proof}
Let $t:=\sqrt{(G^2+Bs^2)/(a^2+B)}$.
At the stationary point, $q_*=(G-at)/B$ and
\eqref{eq:stationary} implies $\sqrt{d(q_*)}=t$.
Thus
\begin{align*}
 z_-(q_*)
 &=\frac{1+a(G-at)/B-t}{2K}\\
 &=\frac{B+aG-(a^2+B)t}{2KB},
\end{align*}
which is \eqref{eq:zminimum}.  Positivity follows independently from
\eqref{eq:vieta}.
\end{proof}

\begin{corollary}[Sharp initial-slope criterion]
Under the baseline assumptions \eqref{eq:base}, the lower real branch
has an interior minimum at a positive $q$ if and only if
\eqref{eq:condition} holds.  Otherwise its minimum is at $q=0$.
\end{corollary}
\begin{proof}
For arbitrary real $G$ satisfying \eqref{eq:base}, the positive
endpoint $q_{\mathrm{coll}}=(G+\sqrt{G^2+Bs^2})/B$ exists.
Equation \eqref{eq:convex} gives strict convexity on its interior,
without requiring $G>0$.  Furthermore,
\[
 \lim_{q\uparrow q_{\mathrm{coll}}}z_-'(q)=+\infty,
\]
because $G-Bq_{\mathrm{coll}}=-\sqrt{G^2+Bs^2}<0$.
Thus an interior stationary point exists precisely when
$z_-'(0)<0$, equivalent by \eqref{eq:initial} to
\eqref{eq:condition}. If the initial derivative is nonnegative,
strict convexity makes $z_-$ increasing for every $q>0$ on the real
interval and its minimum is at zero.
\end{proof}

\section{Interpretation and limitations}
\begin{remark}[Wavenumber versus squared wavenumber]
The strict-convexity statement is about $q=|\mathbf{k}|^2$.
For $k=|\mathbf{k}|>0$, $z_-(k^2)$ has its unique positive
stationary point at $k_*=\sqrt{q_*}$. At that point,
\[
 \frac{d^2}{dk^2}z_-(k^2)\bigg|_{k=k_*}
 =4k_*^2z_-''(q_*)>0.
\]
The even dependence on $k$ also makes $k=0$ stationary when the
function is viewed on the entire signed $k$ axis; it is not the
minimizer under the strict criterion.
\end{remark}

\begin{remark}[No full-theory stability conclusion]
The theorem identifies an exact oscillatory-frequency minimum on the
real-pole interval; it does not eliminate ghost degrees of freedom,
establish a healthy Hamiltonian, or prove nonlinear stability.
Indeed, $d(q)<0$ for $q>q_{\mathrm{coll}}$, so the two $z$ roots
are non-real there.  The corresponding time-frequency roots include
exponentially growing solutions when the untruncated dispersion law is
continued into that regime. The selected scale also depends on
$K,T,G,m^2$; it is not independently calibrated by the algebra alone.
\end{remark}

\noindent\textbf{Related derivation:}
\url{https://github.com/rickyjreyes/geometry_of_resonance/blob/main/docs/cross_paper/FINITE_SCALE_SELECTION.md}.

\end{document}
